\documentclass[main.tex]{subfiles}


\begin{document}

%from pagenumber to up
% \phantomsection 
% \hypertarget{toc}{}
 

 

% номер секции вручную
\setcounter{section}{32
}
\addtocounter{section}{-1} 

\section{Магнитное поле прямого провода с током}


Пусть имеется прямой провод, через который протекает электрический ток. На рис.\,\ref{fig:p147} показана картина линий магнитного поля вокруг провода.


\makeatletter

\pgfkeys{/pgf/decoration/.cd,
         start color/.store in =\startcolor,
         end color/.store in   =\endcolor
}

\pgfdeclaredecoration{width and color change}{initial}{
 \state{initial}[width=0pt, next state=line, persistent precomputation={%
   \pgfmathdivide{50}{\pgfdecoratedpathlength}%
   \let\increment=\pgfmathresult%
   \def\x{0}%
 }]{}
 \state{line}[width=.5pt,   persistent postcomputation={%
     \pgfmathadd@{\x}{\increment}%
     \let\x=\pgfmathresult%
   }]{%
   \pgfsetlinewidth{\x/150*0.035pt+\pgflinewidth}%
   \pgfsetarrows{-}%
   \pgfpathmoveto{\pgfpointorigin}%
   \pgfpathlineto{\pgfqpoint{.75pt}{0pt}}%
   \pgfsetstrokecolor{\endcolor!\x!\startcolor}%
   \pgfusepath{stroke}%
 }
 \state{final}{%
   \pgfsetlinewidth{\pgflinewidth}%
   \pgfpathmoveto{\pgfpointorigin}%
   \color{\endcolor!\x!\startcolor}%
   \pgfusepath{stroke}% 
 }
}

\makeatother


    \begin{figure}[H]
    \centering

    \begin{tikzpicture}[font=\small,            line cap=round,                             >={Stealth[bend,inset=0pt,round,length=3mm, angle'=20]},vec/.style={>={Stealth[inset=0pt,round,length=3.5mm, angle'=20]}}]


\draw[thick,brown] (-1.5,0) coordinate (I) -- (6,0);


%lines
% \foreach \i in {1,2,...,89}{
% \draw[] (0,1) 
% arc [start angle={90*1cm-\i*1cm}, end angle={270*1cm-\i*1cm}, x radius=.5cm, y radius=1cm] 
% % arc [start angle=270, end angle=450, x radius=.5cm, y radius=1cm] 
% ;
% }

\def\xr{.25}
\begin{scope}[]
    
\foreach \i in {0,1,2,3,4}{
\draw[xshift={\i*1cm},line width=.4pt, decoration={width and color change,start color=gray, end color=gray}, decorate] (0.5,{.7pt})               arc [start angle={2}, end angle={179}, x radius={\xr*1cm}, y radius=1cm] 
% arc [start angle=270, end angle=180, x radius=.5cm, y radius=1cm] 
;

\draw[xshift={\i*1cm},line width=.4pt, decoration={width and color change,start color=gray, end color=gray}, decorate] (0.5,{-.7pt})
% arc [start angle={90}, end angle={180}, x radius=.5cm, y radius=1cm] 
arc [start angle=358, end angle=181, x radius={\xr*1cm}, y radius=1cm] 
;

%arrows
\draw[gray,xshift={\i*1cm},->,thin,>={Stealth[bend,inset=0pt,round,length=3.5mm, angle'=20]}] (0.25,{1cm-.5*\pgflinewidth})
% arc [start angle={90}, end angle={180}, x radius=.5cm, y radius=1cm] 
arc [start angle=90, end angle=150, x radius={\xr*1cm+.4\pgflinewidth}, y radius=1cm] 
;


}

\end{scope}

% %arrows
% \draw[red,xshift={1cm},->] (0.25,{.7pt+1cm})
% % arc [start angle={90}, end angle={180}, x radius=.5cm, y radius=1cm] 
% arc [start angle=90, end angle=180, x radius={\xr*1cm}, y radius=1cm] 
% ;


%I
\draw[->,red] ([shift={(0,.2)}]I) --++ (1,0) node[midway, above, black] {$I$}; 

% \draw[red] (0,0) --++ (-1.5,0);
% \draw[red] (6,0) --++ (-1.5,0);

    
    \end{tikzpicture}
\caption{Линии магнитного поля прямого провода с током}
\label{fig:p147}
\end{figure}


Линии поля прямого провода с током являются окружностями. Центры этих окружностей лежат на проводе, при этом окружности лежат в плоскостях, перпендикулярных проводу (см.~рис.\,\ref{fig:p147}).

% Удобно использовать термин <<ближайшая сторона линии>>. Так, по ближайшим сторонам линий на рис.\,\ref{fig:p147}  

Через любую точку пространства вне провода можно провести соответствующую линию (в форме окружности) его магнитного поля. С удалением от провода густота линий поля убывает (убывает индукция магнитного поля).

Для определения направления линий магнитного поля прямого тока можно пользоваться следующим правилом.
%
\begin{law}
    \textbf{Правило правой руки.} Если большой палец правой руки расположить так, чтобы он указывал направление тока в проводе, то остальные согнутые пальцы (охватывающие провод) укажут направление линий магнитного поля, создаваемого этим током.
\end{law}
%
Рисунок~\ref{fig:p148} иллюстрирует применение этого правила.  

    \begin{figure}[H]
    \centering
\includegraphics[]{el/36wire_pic.pdf}%
\caption{Правило правой руки}
\label{fig:p148}
\end{figure}


Сопоставляя рисунки~\ref{fig:p147} и~\ref{fig:p148}, можно видеть, что согнутые пальцы на рис.\,\ref{fig:p148} указывают направление линий магнитного поля вокруг провода.

\medskip

\noindent\textbf{Задача.} Два одинаковых изолированных тонких провода с противоположно направленными равными токами плотно вставлены в узкую пластиковую трубку. Что можно сказать о магнитном поле вокруг трубки? 

\smallskip

\noindent\emph{Решение.} 
% Одинаковые токи характеризуются одинаковыми картинами линий магнитного поля. 
С помощью правила правой руки можно убедиться, что поля токов компенсируют друг друга~--- магнитное поле вокруг трубки отсутствует.















 
\end{document}
